\section{Validacion RealBike con datos reales y OSRM}
\label{sec:realbike-final}

RealBike se plantea como una validacion externa de ADR-EDA sobre datos reales de bike-sharing station-based. A diferencia de los benchmarks sinteticos, aqui las estaciones, coordenadas y viajes proceden de Citi Bike y Divvy. La decision optimizada por ADR-EDA es un vector de inventario objetivo por estacion; la rutina de recolocacion traduce ese vector en movimientos de bicicletas, y la demanda observada actualiza el estado del sistema.

La motivacion es operacional: reoptimizar cada hora puede mejorar el servicio, pero tambien incrementa los kilometros de recolocacion y el coste computacional. ADR-EDA introduce una puerta adaptativa que decide cuando actualizar la distribucion aprendida y cuando reutilizar la solucion heredada.

\begin{table}[htbp]
\centering
\scriptsize
\caption{Políticas comparadas en RealBike}
\label{tab:realbike-politicas}
\begin{tabularx}{\textwidth}{
    >{\raggedright\arraybackslash}p{2.7cm}
    >{\raggedright\arraybackslash}p{2.3cm}
    >{\centering\arraybackslash}p{1.5cm}
    >{\centering\arraybackslash}p{1.3cm}
    >{\raggedright\arraybackslash}X
}
\toprule
Política & Reoptimiza & Reutiliza & Usa EDA & Rol \\
\midrule
Estática 2/día & 03:00 y 15:00 & no & no & Referencia periódica inspirada en reequilibrado estático\\
Reequilibrado horario & cada hora & no & no & Referencia determinista que recalcula un objetivo en todos los periodos\\
Horizonte rodante & cada hora & no & no & Referencia operativa que usa una predicción de demanda de corto plazo\\
EDA horaria & cada hora & no & sí & Referencia intensiva con EDA sin reutilización adaptativa\\
ADR-EDA & adaptativo & sí & sí & Política propuesta: actualiza o reutiliza según las señales del dominio\\
\bottomrule
\end{tabularx}
\end{table}


\begin{table}[htbp]
\centering
\scriptsize
\caption{Resultados principales Citi Bike con matriz OSRM, 336 horas y 10 semillas}
\label{tab:realbike-citibike-osrm}
\resizebox{\textwidth}{!}{%
\begin{tabular}{lllllll}
\hline
Política & Utilidad & Servicio & Km & Reopt& Reutil& Rechazos \\
\hline
sin reequilibrado & 4766.8 & 0.8640 & 0.0 & 0.0 & 0.0 & 1237.2 \\
estática 2/día & 8115.3 & 0.9743 & 2037.5 & 23.0 & 0.0 & 234.0 \\
reeq. horario & 8552.4 & 0.9931 & 4107.5 & 336.0 & 0.0 & 63.0 \\
horizonte rodante & 8628.8 & 0.9953 & 4012.9 & 336.0 & 0.0 & 43.0 \\
selectiva & 8591.3 & 0.9895 & 2040.3 & 152.0 & 0.0 & 95.7 \\
ADR aproximado & 8435.3 & 0.9843 & 1903.8 & 87.5 & 61.9 & 142.8 \\
EDA horaria & 8682.5 & 0.9951 & 2809.4 & 336.0 & 0.0 & 44.2 \\
ADR-EDA & 8582.6 & 0.9902 & 2227.4 & 198.6 & 26.1 & 88.8 \\
\hline
\end{tabular}
}
\end{table}


En la campana Citi Bike con matriz OSRM, ADR-EDA supera claramente al baseline estatico dos veces al dia y se aproxima a las politicas dinamicas agresivas. Frente a rolling horizon reduce kilometros y reoptimizaciones, aunque con una ligera perdida de servicio. Frente a EDA horaria sin reutilizacion, ADR-EDA mantiene un nivel de servicio alto con menos actualizaciones, lo que justifica la puerta adaptativa.

\begin{table}[htbp]
\centering
\scriptsize
\caption{Comparación pareada de ADR-EDA frente a referencias Citi Bike OSRM}
\label{tab:realbike-citibike-pairwise}
\resizebox{\textwidth}{!}{%
\begin{tabular}{lllll}
\hline
Comparación & Delta utilidad & Delta servicio & Delta km & Delta reopt\\
\hline
ADR-EDA frente a estática 2/día & 467.3 & 0.0160 & 189.9 & 175.6 \\
ADR-EDA frente a reeq. horario & 30.2 & -0.0028 & -1880.1 & -137.4 \\
ADR-EDA frente a horizonte rodante & -46.2 & -0.0050 & -1785.5 & -137.4 \\
ADR-EDA frente a EDA horaria & -99.9 & -0.0049 & -582.0 & -137.4 \\
\hline
\end{tabular}
}
\end{table}


Adicionalmente, siguiendo la literatura reciente sobre reequilibrado dinamico, se ejecuta una validacion cronologica evento-a-evento. Esta evaluacion evita depender exclusivamente de demanda agregada por hora: cada alquiler y devolucion se aplica en su instante observado, mientras que la politica sigue decidiendo en fronteras horarias.

\begin{table}[htbp]
\centering
\footnotesize
\caption{Resultados evento-a-evento con Citi Bike y Divvy, 168 horas, 10 semillas y OSRM}
\label{tab:realbike-event-replay}
\begin{tabular}{llrrrrrr}
\hline
Sistema & Criterio & Servicio & Utilidad & $\Delta U$ horiz& Km & Reopt& Km ahorr. horiz\\
\hline
Citi Bike & 1/4 sensible & 1.0000 & 4036.2 & 5.6 & 1478.5 & 167.4 & 440.8 \\
Citi Bike & 2/8 reutilizador & 0.9997 & 4024.1 & -6.5 & 1558.9 & 84.0 & 360.5 \\
Divvy & 1/4 sensible & 0.9804 & 1913.6 & 52.9 & 1164.1 & 168.0 & 920.7 \\
Divvy & 2/8 reutilizador & 0.9760 & 1861.3 & 0.6 & 1319.7 & 89.9 & 765.2 \\
\hline
\end{tabular}
\end{table}


Los resultados evento-a-evento muestran que ADR-EDA no colapsa al pasar a una simulacion mas fina. En Citi Bike, con gate agresivo, iguala el servicio de rolling horizon y reduce kilometros. En Divvy, el mismo gate mejora la utilidad frente a rolling horizon y tambien reduce kilometros. El gate mas reutilizador reduce aproximadamente a la mitad las reoptimizaciones, a costa de una perdida moderada, lo que ofrece una palanca interpretable para ajustar la politica al coste operativo.

La conclusion defendible es que ADR-EDA funciona como politica coste-consciente para reequilibrado dinamico station-based: aprende distribuciones de inventario objetivo y decide cuando actualizarlas o reutilizarlas. Este bloque no pretende reproducir todos los detalles industriales de una empresa de bike-sharing, pero si aporta una validacion realista, reproducible y alineada con literatura mediante datos reales, distancias OSRM, baselines operativos y replay cronologico de eventos.
